% !TeX root = ../../Book4.tex
% 纲要标题：## 第11章　同调维数与零调对象
% 纲要标签：b4:ch:10
% 纲要位置：计划纲要.md，第 3463 行。

\chapter{同调维数与零调对象}
\label{b4:ch:10}
\BookChapterNumber{b4:ch:10}

\begin{chapterabstract}
本章以消解长度和高次 Ext、Tor 的消失刻画同调维数，研究何时可截短消解，以及用零调对象计算导出函子的条件。最后研究复合函子的导出，区分正合性与内射像的零调条件在两种直接复合公式中的作用，并引出 Grothendieck 谱序列。
\end{chapterabstract}

\begin{learninggoals}
\begin{itemize}
\item 区分投射、内射、平坦维数及环的左、右整体维数。
\item 用 Ext、Tor 的消失判据证明消解可以截短，而不把某一消解的表面长度当成不变量。
\item 判断指定对象是否对指定函子零调，并用适当消解计算导出函子。
\item 分清复合中前函子正合、后函子正合及内射像零调三个条件的不同作用。
\end{itemize}
\end{learninggoals}

令 \(R=k[\varepsilon]/(\varepsilon^2)\)，其中 \(k\) 是域。商模 \(k=R/(\varepsilon)\) 由 \(R\twoheadrightarrow k\) 给出；核是 \(\varepsilon R\)，它又是乘 \(\varepsilon\) 的像，而乘 \(\varepsilon\) 的核仍是 \(\varepsilon R\)。因此沿着乘 \(\varepsilon\) 可以不断向左接出自由模。单凭这条无限消解，还不能排除另有一条有限投射消解；关键是对它取 \(\operatorname{Hom}_R(-,k)\) 后微分全为零，得到每个正次数的 \(\operatorname{Ext}_R^i(k,k)\) 都非零。有限投射消解必使足够高次的 Ext 消失，所以这里确实无法用有限投射消解替代。这与整数上的 \(\mathbb Z/n\mathbb Z\) 有长度一的自由消解形成对照。本章将说明消解长度怎样由 Ext、Tor 的高次消失检测，以及何时能换用对某个函子零调的对象。

\ProseParagraphBreak
本章直接使用\chapref{b4:ch:07}、\chapref{b4:ch:08}的消解、导出函子长正合列以及\secref{b4:sec:0805}建立的维数平移公式；\secref{b4:sec:1002}用这些公式判断何时能够截短消解，并估计同调维数。局部环上有限生成模的例子调用第三卷第17.1—17.2节已证明的极小自由消解和投射维数判据。\secref{b4:sec:1004}列出 Grothendieck 谱序列的适用条件，证明见\chapref{b4:ch:15}。

\ProseParagraphBreak
同调维数可对照 Weibel 的《An Introduction to Homological Algebra》\cite[Section 4.1]{Weibel1994HomologicalAlgebra}；零调对象与消解计算见\cite[F-Acyclic Objects 2.4.3, Exercise 2.4.3, Section 2.5]{Weibel1994HomologicalAlgebra}。这里按本书的左模约定书写 Ext，并在每个 Tor 判据中明确另一变量为右模。

% !TeX root = ../../Book4.tex
% 纲要标题：### 11.1 同调维数
% 纲要标签：b4:sec:1001
% 纲要位置：计划纲要.md，第 3465 行。

\section{同调维数}
\label{b4:sec:1001}

一个模的某条消解可能很长，但这未必来自模本身。我们要问的是它能否有更短的消解，以及导出群怎样检测消解能在何处结束。

% 纲要内容：
%
% * 投射、内射及平坦维数
% * 整体维数
% * 有限维数与消失判据
%

\subsection{投射、内射及平坦维数}
\label{b4:subsec:1001:01}

具体消解的长度按其最高非零次数计数，不把增广目标计作消解的一项。例如，非零投射模 \(P\) 有长度零的消解 \(0\to P\xrightarrow{1_P}P\to0\)。取非零投射模 \(Q\)，在任意两个相邻的正次数上附加可缩复形
\[
0\longrightarrow Q\xrightarrow{1_Q}Q\longrightarrow0,
\]
仍得到 \(P\) 的投射消解，却可使最高非零次数任意大。去掉这部分是对具体消解作同伦化简；要定义模本身的维数，则须在全部允许的消解中取能够达到的最短长度，而不能只看眼前一条消解有多少非零项。

\begin{definition}\label{b4:def:homological-dimensions}
设 \(R\) 为含幺环，\(M\ne0\) 为幺性左 \(R\) 模。对非负整数 \(n\)，考虑正合列
\[
 0\longrightarrow P_n\longrightarrow\cdots\longrightarrow P_0
 \longrightarrow M\longrightarrow0
\]
若每个 \(P_i\)（\(0\le i\le n\)）投射，则在这类列中取最小长度 \(n\)，称为 \(M\) 的\term[toushe-weishu]{投射维数}[projective dimension]，记为 \(\operatorname{pd}_RM\)。把投射模换为平坦左 \(R\) 模，得到\term[pingtan-weishu]{平坦维数}[flat dimension] \(\operatorname{fd}_RM\)。若改用正合列
\[
 0\longrightarrow M\longrightarrow I^0\longrightarrow\cdots\longrightarrow I^n\longrightarrow0
\]
且每个 \(I^j\)（\(0\le j\le n\)）内射，则最小长度称为\term[neishe-weishu]{内射维数}[injective dimension]，记为 \(\operatorname{id}_RM\)。没有有限长度消解时，对应维数为 \(\infty\)；对零模三者均约定为 \(-\infty\)。右模的维数定义在 \(R^{\mathrm{op}}\) 上。零模的约定不表示存在负无限长度的消解，而是使零项不影响维数比较中的最大值，例如 \(M\oplus0=M\)，相应的最大值也应保持不变。本章对有限整数 \(r\) 约定
\[
\begin{gathered}
 (-\infty)\pm r=-\infty,\qquad (+\infty)\pm r=+\infty,
 \\
 \max\{-\infty,d\}=d
 \quad\bigl(d\in\mathbb Z\cup\{\pm\infty\}\bigr),
\end{gathered}
\]
并约定空集的上确界为 \(\sup\varnothing=-\infty\)。
\end{definition}
\symidx{\(\operatorname{pd}_RM\)：模的投射维数}
\symidx{\(\operatorname{id}_RM\)：模的内射维数}
\symidx{\(\operatorname{fd}_RM\)：模的平坦维数}

这些方向来自\secref{b4:sec:0701}中的构造：投射消解从到 \(M\) 的满射开始，逐次对核取投射满射，向左延伸；内射消解从 \(M\) 到内射模的单射开始，逐次把余核嵌入内射模，向右延伸。平坦消解与投射消解同向，只减弱了各项的要求。长度零时，正合列中唯一的消解项与 \(M\) 同构，所以非零模的三种维数为零，分别等价于它本身投射、内射、平坦。每个投射模都平坦，故每条投射消解也是平坦消解，总有 \(\operatorname{fd}_RM\le\operatorname{pd}_RM\)。

\ProseParagraphBreak
这个不等式可能严格。令 \(S=\mathbb Z\setminus\{0\}\)，则 \(\mathbb Q=S^{-1}\mathbb Z\)。第二卷命题6.4.6、6.4.7给出模的局部化与张量积的自然同构，以及局部化的正合性，利用整数环的交换性得到
\[
A\otimes_{\mathbb Z}\mathbb Q\cong S^{-1}A
\]
对交换群 \(A\) 自然，\(-\otimes_{\mathbb Z}\mathbb Q\) 保持短正合列，故 \(\mathbb Q\) 平坦。它却不投射：若投射，第二卷定理7.1.2使它成为某个自由交换群的直和因子，而\cref{b4:thm:pid-free-submodule}使这个子群自由。非零自由交换群中任一基向量都不能被 \(2\) 整除，与 \(\mathbb Q\) 可除矛盾。因此 \(\operatorname{pd}_{\mathbb Z}\mathbb Q>0\)。另一方面，\(\mathbb Z\) 是主理想整环，同一定理给出 \(\mathbb Q\) 的长度至多一的自由消解，故其投射维数不超过一。两者合得
\[
\operatorname{fd}_{\mathbb Z}\mathbb Q=0,
\qquad \operatorname{pd}_{\mathbb Z}\mathbb Q=1.
\]

\ProseParagraphBreak
同一个交换群改变底环后也会改变维数。设 \(p\) 为素数，\(\mathbb Z/p\mathbb Z\) 作为 \(\mathbb F_p\)-模投射维数为零，作为整数模的投射维数为一，而作为 \(\mathbb Z/p^2\mathbb Z\)-模有无限投射维数。对最后一个底环，令 \(S=\mathbb Z/p^2\mathbb Z\)，乘 \(p\) 的核与像均为 \(pS\)，反复使用这个映射便得到自由消解 \(\cdots\xrightarrow{p}S\xrightarrow{p}S\to\mathbb Z/p\mathbb Z\to0\)。施加 \(\operatorname{Hom}_S(-,\mathbb Z/p\mathbb Z)\) 后微分为零，故任意正次数的 Ext 自群都是非零的 \(\mathbb Z/p\mathbb Z\)，与有限投射消解在高次数没有项矛盾。\cref{b4:ex:10-three-base-rings}进一步比较这三个底环下的完整计算；群的大小不能决定这里的差别。

\subsection{整体维数}
\label{b4:subsec:1001:02}

\begin{definition}\label{b4:def:global-dimension}
设 \(R\) 为非零含幺环。所有幺性左 \(R\) 模投射维数的上确界
\[
 \operatorname{l.gl.dim}R=\sup\{\,\operatorname{pd}_RM:M\;\text{为左}\;R\text{-模}\,\}
\]
称为 \(R\) 的\term[zuozhengtiweishu]{左整体维数}[left global dimension]。右整体维数定义为 \(\operatorname{r.gl.dim}R=\operatorname{l.gl.dim}R^{\mathrm{op}}\)。交换环时记为 \(\operatorname{gl.dim}R\)。
\end{definition}
\symidx{\(\operatorname{l.gl.dim}R,\ \operatorname{r.gl.dim}R\)：环的左整体维数与右整体维数}
\symidx{\(\operatorname{gl.dim}R\)：交换环的整体维数}

域及除环的所有模都有基，故相应整体维数为零。非交换环的左、右整体维数不能不加说明地混写。这里的上确界也包括无限生成模；只讨论有限生成模时，得到的是另一个受限制的问题。

\ProseParagraphBreak
单个左模 \(M\) 投射，只保证以它为目标的满射有截面；左整体维数为零则要求每个左模都投射，因而每条短正合列 \(0\to A\to B\to C\to0\) 的末项 \(C\) 都投射，使该列分裂。反过来，若整个左模范畴的短正合列都分裂，对任意左模 \(M\) 取自由满射 \(F\to M\)，所得短正合列分裂，使 \(M\) 成为自由模的直和因子，因而投射。这就证明左整体维数为零当且仅当每条短正合左模列分裂；它测量的是整个模范畴的性质。第二卷第十三章的半单模结构理论进一步把这种性质与半单环联系起来。

\subsection{有限维数与消失判据}
\label{b4:subsec:1001:03}

消解能在第 \(d\) 次结束，会使所有更高次导出群消失。反过来，并不需要逐个检验这些高次数：维数平移把第 \(d+1\) 次的消失化为第 \(d\) 个合冲模或余合冲模的一次消失，再用投射、内射、平坦的一次判据，就能使消解在这一位置结束。这里必须让测试模遍历相应的全部模。

\begin{theorem}\label{b4:thm:dimension-vanishing}
设 \(R\) 为含幺环，\(M\) 为幺性左 \(R\) 模，\(d\ge0\) 为整数。有以下等价条件：
\begin{enumerate}
\item \(\operatorname{pd}_RM\le d\)，当且仅当对每个幺性左 \(R\) 模 \(N\) 有 \(\operatorname{Ext}_R^{d+1}(M,N)=0\)；这又等价于所有 \(i>d\) 的这些 Ext 群为零。
\item \(\operatorname{id}_RM\le d\)，当且仅当对每个幺性左 \(R\) 模 \(N\) 有 \(\operatorname{Ext}_R^{d+1}(N,M)=0\)；这又等价于所有 \(i>d\) 的这些 Ext 群为零。
\item \(\operatorname{fd}_RM\le d\)，当且仅当对每个幺性右 \(R\) 模 \(T\) 有 \(\operatorname{Tor}_{d+1}^R(T,M)=0\)；这又等价于所有 \(i>d\) 的这些 Tor 群为零。
\end{enumerate}
\end{theorem}
\begin{proof}
投射情形，长度至多为 \(d\) 的投射消解在第 \(d+1\) 次及以上没有项，故施加 \(\operatorname{Hom}_R(-,N)\) 后，高于 \(d\) 次的 Ext 为零。反过来，取任意投射消解，令 \(K_d=\Omega^dM\)，其中 \(K_0=M\)。由\cref{b4:prop:iterated-dimension-shift}，对每个左模 \(N\) 都有
\[
\operatorname{Ext}_R^1(K_d,N)\cong\operatorname{Ext}_R^{d+1}(M,N).
\]
若右边对所有 \(N\) 消失，则\cref{b4:thm:ext-projective-injective}使 \(K_d\) 投射。\(d>0\) 时，截取的正合列为
\[
0\longrightarrow K_d\longrightarrow P_{d-1}\longrightarrow\cdots
\longrightarrow P_0\longrightarrow M\longrightarrow0.
\]
把 \(K_d\) 放在第 \(d\) 次，就得到长度至多为 \(d\) 的投射消解；\(d=0\) 时判据直接使 \(M\) 投射。这也使所有更高次 Ext 自动消失。

内射情形，有限内射消解直接给出全部更高次 Ext 的消失。反过来，固定内射消解 \(M\to I^\bullet\)，令 \(L^0=M\)，并逐次取
\[
L^{j+1}=\operatorname{coker}(L^j\hookrightarrow I^j).
\]
其中每个单射由消解的正合性给出，见\cref{b4:prop:iterated-dimension-shift}中的余合冲构造。同一命题给出
\[
\operatorname{Ext}_R^1(N,L^d)\cong\operatorname{Ext}_R^{d+1}(N,M).
\]
若右边对所有左模 \(N\) 消失，一次内射判据使 \(L^d\) 内射。\(d>0\) 时以它作为最后一项，得到
\[
0\longrightarrow M\longrightarrow I^0\longrightarrow\cdots
\longrightarrow I^{d-1}\longrightarrow L^d\longrightarrow0,
\]
这是长度至多为 \(d\) 的内射消解；\(d=0\) 时检测的对象就是 \(L^0=M\)。

平坦情形，\cref{b4:thm:tor-flat-criterion}给出平坦左模与任意右模的正次 Tor 全部为零。把长度至多为 \(d\) 的平坦消解拆成短正合列，由\cref{b4:thm:tor-long-exact}逐次平移：高于 \(d\) 次的 Tor 最终化为最左平坦项的正次 Tor，故为零。反过来，可先取一个自由消解，仍令 \(K_d=\Omega^dM\)、\(K_0=M\)。\cref{b4:prop:iterated-dimension-shift}给出
\[
\operatorname{Tor}_1^R(T,K_d)\cong\operatorname{Tor}_{d+1}^R(T,M)
\]
对每个右模 \(T\) 成立。一次 Tor 对所有右模消失，使 \(K_d\) 平坦；\(d>0\) 时用与投射情形相同的截短列，原自由项与最左的 \(K_d\) 均平坦，因而得到长度至多为 \(d\) 的平坦消解。\(d=0\) 时则直接检测 \(M\) 的平坦性。这里取的仍是向左的合冲模，与内射情形的余合冲模不同。
\end{proof}

\begin{corollary}\label{b4:cor:global-dimension-ext}
设 \(R\) 为非零含幺环。左整体维数也等于所有幺性左 \(R\) 模内射维数的上确界，并且等于
\[
 \sup\{\,n\ge0:\operatorname{Ext}_R^n(M,N)\ne0
       \;\text{对某些左}\;R\text{-模}\;M,N\,\}.
\]
\end{corollary}
\begin{proof}
对每个整数 \(d\ge0\)，由上一条定理得到
\[
\begin{aligned}
\bigl(\forall M,\ \operatorname{pd}_RM\le d\bigr)
&\Longleftrightarrow
\bigl(\forall M\,\forall N,\ \operatorname{Ext}_R^{d+1}(M,N)=0\bigr)\\
&\Longleftrightarrow
\bigl(\forall N,\ \operatorname{id}_RN\le d\bigr),
\end{aligned}
\]
其中 \(M,N\) 都遍历幺性左 \(R\) 模。因而两种维数的全体取值具有相同的有限上界；没有有限上界时，两边的上确界都为 \(\infty\)，且对每个 \(d\ge0\) 都有一对 \(M,N\) 使第 \(d+1\) 次 Ext 非零。若共同上确界是有限整数 \(g>0\)，在 \(d=g-1\) 时上述统一消失不能成立，故有第 \(g\) 次的非零 Ext，而高于 \(g\) 次的 Ext 全为零。\(R\ne0\) 保证 \(\operatorname{Hom}_R(R,R)\ne0\)，涵盖 \(g=0\) 的情形。这证明两个上确界及所列 Ext 次数的上确界相等，并不要求同一个模的投射维数与内射维数相等。
\end{proof}

因此非域主理想整环 \(R\) 的整体维数恰为一：\cref{b4:thm:pid-free-submodule}给出所有模长度至多一的自由消解；取非零非单位 \(a\)，由\cref{b4:exm:ext-cyclic}得 \(\operatorname{Ext}_R^1(R/(a),R)=R/(a)\ne0\)，给出下界。

\begin{example}\label{b4:exm:infinite-dimension-dual-numbers}
设 \(k\) 为域，\(R=k[\varepsilon]/(\varepsilon^2)\)。把 \(k\) 识别为商模 \(R/(\varepsilon)\)，并由 \(R\) 的交换性同时视为左模与右模。它的周期自由消解为
\[
 \cdots\xrightarrow{\varepsilon}R\xrightarrow{\varepsilon}R
 \xrightarrow{\varepsilon}R\longrightarrow k\longrightarrow0.
\]
乘 \(\varepsilon\) 的核与像均为 \(\varepsilon R\)，也等于商映射的核，所以该列正合。它无限长，还不能排除 \(k\) 另有有限长度的消解。对它施加 \(\operatorname{Hom}_R(-,k)\) 或 \(-\otimes_Rk\)，因为 \(\varepsilon\) 在 \(k\) 上作用为零，全部微分都为零，而每一项都同构于 \(k\)。因此
\[
\operatorname{Ext}_R^i(k,k)\cong k,
\qquad \operatorname{Tor}_i^R(k,k)\cong k
\quad(i\ge0).
\]
对任意有限 \(d\ge0\)，第 \(d+1\) 次仍非零，\cref{b4:thm:dimension-vanishing}的三个判据遂分别排除有限投射、内射、平坦消解。于是 \(\operatorname{pd}_Rk=\operatorname{id}_Rk=\operatorname{fd}_Rk=\infty\)，且 \(\operatorname{gl.dim}R=\infty\)。无限维数的依据是任意高次仍有非零导出群，与所选消解的表面长度不同。
\end{example}

维数与消失条件的系统对应可参阅 Weibel\cite[Section 4.1, pd Lemma 4.1.6, id Lemma 4.1.8, fd Lemma 4.1.10]{Weibel1994HomologicalAlgebra}。

% !TeX root = ../../Book4.tex
% 纲要标题：### 11.2 截短消解与同调维数估计
% 纲要标签：b4:sec:1002
% 纲要位置：计划纲要.md，第 3471 行。

\section{截短消解与同调维数估计}
\label{b4:sec:1002}

高次导出群是否消失，可以转化为消解末端的合冲模是否投射。这个判据既判断给定消解能否截短，也说明局部环上的投射维数可以怎样检测。

% 纲要内容：
%
% * 截短消解及合冲模
% * 短正合列中的同调维数不等式
% * 第三卷局部投射维数所需的最低平移版本已在第9.5节证明；本节作一般化整理
% * 局部投射维数的讨论作为交换代数应用，可在学习第三卷后回读；不参与前两小节的证明
%

\subsection{截短消解及合冲模}
\label{b4:subsec:1002:01}

维数的定义只要求存在一条足够短的消解；下面的判据则作用于任意一条已经给定的消解。若 \(\operatorname{pd}_RM\le d\)，无论前 \(d\) 步选用了哪些投射满射，得到的第 \(d\) 个合冲模都投射，因而都能在此截短。内射方向有同样的结论；对于平坦维数，当前给定的各项只需平坦，不必投射。

\begin{proposition}\label{b4:prop:truncation-dimension}
设 \(R\) 为含幺环，\(M\) 为幺性左 \(R\) 模，\(d\ge1\) 为整数。给定幺性左 \(R\) 模正合列
\[
 0\longrightarrow K_d\longrightarrow P_{d-1}\longrightarrow\cdots
 \longrightarrow P_0\longrightarrow M\longrightarrow0,
\]
其中各 \(P_i\) 投射。则 \(\operatorname{pd}_RM\le d\) 当且仅当 \(K_d\) 投射。若幺性左 \(R\) 模序列 \(0\to M\to I^0\to\cdots\to I^{d-1}\to L^d\to0\) 正合且各 \(I^i\) 内射，则 \(\operatorname{id}_RM\le d\) 当且仅当 \(L^d\) 内射。若第一列的 \(P_i\) 仅要求平坦，则 \(\operatorname{fd}_RM\le d\) 当且仅当 \(K_d\) 平坦。
\end{proposition}
\begin{proof}
若末端对象具有所需性质，将它作为第 \(d\) 个投射项、内射项或平坦项，给定列本身就是长度至多 \(d\) 的相应消解。

反向令 \(K_0=M\)，把第一列拆成 \(0\to K_{j+1}\to P_j\to K_j\to0\)（\(0\le j<d\)）；这里 \(P_j\to K_j\) 是原微分到其像的满射，最后的核就是给定的 \(K_d\)。若各 \(P_j\) 投射，对任意幺性左模 \(N\) 依次用\cref{b4:thm:dimension-shift-basic}，得到
\[
 \operatorname{Ext}_R^1(K_d,N)\simeq\operatorname{Ext}_R^{d+1}(M,N).
\]
\(\operatorname{pd}_RM\le d\) 使右侧对所有 \(N\) 消失；\cref{b4:thm:ext-projective-injective}的一次判据便使这个给定的 \(K_d\) 投射。

内射列则令 \(L^0=M\)，依次取余核，拆成 \(0\to L^j\to I^j\to L^{j+1}\to0\)（\(0\le j<d\)）。第二变量的平移\cref{b4:prop:injective-dimension-shift}给出
\[
 \operatorname{Ext}_R^1(N,L^d)\simeq\operatorname{Ext}_R^{d+1}(N,M).
\]
若 \(\operatorname{id}_RM\le d\)，右侧对所有幺性左模 \(N\) 消失，一次 Ext 判据使 \(L^d\) 内射。

平坦版本仍拆第一列，但不使用 \(P_j\) 的投射性。对任意幺性右模 \(T\)，\cref{b4:thm:tor-flat-criterion}使 \(\operatorname{Tor}_i^R(T,P_j)=0\) 对每个 \(i>0\) 成立；\cref{b4:thm:tor-long-exact}于是逐步给出
\[
 \operatorname{Tor}_1^R(T,K_d)\simeq\operatorname{Tor}_{d+1}^R(T,M).
\]
若 \(\operatorname{fd}_RM\le d\)，右侧由\cref{b4:thm:dimension-vanishing}对所有 \(T\) 消失，同一平坦判据使 \(K_d\) 平坦。
\end{proof}

因此某个 \(d\) 步合冲模投射时，任何投射消解的 \(d\) 步合冲模都投射。这些合冲模不必同构，但加入投射直和项后同构，即稳定同构。具体地，两个投射消解的第一步给出短正合列 \(0\to K_1\to P_0\to M\to0\) 与 \(0\to K'_1\to P'_0\to M\to0\)，由\cref{b4:lem:schanuel}有
\[
 K_1\oplus P'_0\cong K'_1\oplus P_0.
\]
对第二步，记原消解给出的满射为 \(\pi_1:P_1\to K_1\)、\(\pi'_1:P'_1\to K'_1\)。分别补入恒等映射，得到
\[
\begin{aligned}
0&\longrightarrow K_2\longrightarrow P_1\oplus P'_0
 \xrightarrow{\pi_1\oplus1}K_1\oplus P'_0\longrightarrow0,\\
0&\longrightarrow K'_2\longrightarrow P'_1\oplus P_0
 \xrightarrow{\pi'_1\oplus1}K'_1\oplus P_0\longrightarrow0.
\end{aligned}
\]
恒等分量的核为零，所以两条新列的核仍分别是 \(K_2,K'_2\)；中间项都是投射模。利用第一步的同构将两个右端识别，再用 Schanuel 引理，得到
\[
 K_2\oplus P'_1\oplus P_0\cong K'_2\oplus P_1\oplus P'_0.
\]
同样，若在第 \(j\) 步已有 \(K_j\oplus Q\cong K'_j\oplus Q'\)，其中 \(Q,Q'\) 投射，就比较 \(P_j\oplus Q\to K_j\oplus Q\) 与 \(P'_j\oplus Q'\to K'_j\oplus Q'\)，恒等分量仍不改变核。再次应用该引理，便得到第 \(j+1\) 步的稳定同构。归纳可得投射模 \(Q,Q'\)，使 \(K_d\oplus Q\cong K'_d\oplus Q'\)。同调维数检测的是这些合冲模何时投射，并不要求它们作为模相互同构。

\subsection{短正合列中的同调维数不等式}
\label{b4:subsec:1002:02}

\begin{proposition}\label{b4:prop:dimension-short-exact}
设 \(R\) 为含幺环，\(0\to A\to B\to C\to0\) 为幺性左 \(R\) 模短正合列。以下所有维数均在 \(R\) 上计算。投射维数满足
\[
\begin{aligned}
 \operatorname{pd}B&\le\max\{\operatorname{pd}A,\operatorname{pd}C\},\\
 \operatorname{pd}A&\le\max\{\operatorname{pd}B,\operatorname{pd}C-1\},\\
 \operatorname{pd}C&\le\max\{\operatorname{pd}B,\operatorname{pd}A+1\}.
\end{aligned}
\]
平坦维数满足相同的三式。内射维数满足
\[
\begin{aligned}
 \operatorname{id}B&\le\max\{\operatorname{id}A,\operatorname{id}C\},\\
 \operatorname{id}A&\le\max\{\operatorname{id}B,\operatorname{id}C+1\},\\
 \operatorname{id}C&\le\max\{\operatorname{id}B,\operatorname{id}A-1\}.
\end{aligned}
\]
这里维数取值及 \(\pm\infty\) 的运算遵循\cref{b4:def:homological-dimensions}中的约定。
\end{proposition}
\begin{proof}
若某一右端为 \(\infty\)，该不等式直接成立。右端为负的情形由短列的正合性迫使被估计的模为零，结论也直接成立。以下考虑非负有限上界；零次 Ext、Tor 分别按 Hom、张量积解释。

固定幺性左模 \(N\)，\cref{b4:thm:ext-long-exact}在第一变量给出五项片段
\[
\begin{aligned}
 \operatorname{Ext}^{i-1}(A,N)&\longrightarrow\operatorname{Ext}^i(C,N)
 \longrightarrow\operatorname{Ext}^i(B,N)\\
 &\longrightarrow\operatorname{Ext}^i(A,N)
 \longrightarrow\operatorname{Ext}^{i+1}(C,N)
 \qquad(i\ge1).
\end{aligned}
\]
估计 \(B\) 时，取 \(i>\max\{\operatorname{pd}A,\operatorname{pd}C\}\)，它两侧的同次数 Ext 均为零，故 \(\operatorname{Ext}^i(B,N)=0\)。估计 \(A\) 时，条件
\(i>\max\{\operatorname{pd}B,\operatorname{pd}C-1\}\)
同时给出 \(i>\operatorname{pd}B\) 与 \(i+1>\operatorname{pd}C\)，使夹住 \(\operatorname{Ext}^i(A,N)\) 的两项消失。估计 \(C\) 时，条件
\(i>\max\{\operatorname{pd}B,\operatorname{pd}A+1\}\)
则给出 \(i>\operatorname{pd}B\) 与 \(i-1>\operatorname{pd}A\)，使 \(\operatorname{Ext}^{i-1}(A,N)\) 与 \(\operatorname{Ext}^i(B,N)\) 消失。这说明三个最大值中的次数偏移恰好来自 \(i+1\) 和 \(i-1\)。以上对每个 \(N\) 成立，\cref{b4:thm:dimension-vanishing}便给出三条投射维数不等式。

平坦情形固定幺性右模 \(T\)，\cref{b4:thm:tor-long-exact}给出
\[
\begin{aligned}
 \operatorname{Tor}_{i+1}(T,C)&\longrightarrow\operatorname{Tor}_i(T,A)
 \longrightarrow\operatorname{Tor}_i(T,B)\\
 &\longrightarrow\operatorname{Tor}_i(T,C)
 \longrightarrow\operatorname{Tor}_{i-1}(T,A)
 \qquad(i\ge1).
\end{aligned}
\]
\(B\) 仍被 \(A,C\) 的同次数项夹住；\(A\) 被 \(C\) 的 \(i+1\) 次与 \(B\) 的 \(i\) 次项夹住；\(C\) 被 \(B\) 的 \(i\) 次与 \(A\) 的 \(i-1\) 次项夹住。消失所需的次数条件与投射情形相同，故得到同样的三式。

内射情形固定幺性左模 \(N\)，改用第二变量的片段
\[
\begin{aligned}
 \operatorname{Ext}^{i-1}(N,C)&\longrightarrow\operatorname{Ext}^i(N,A)
 \longrightarrow\operatorname{Ext}^i(N,B)\\
 &\longrightarrow\operatorname{Ext}^i(N,C)
 \longrightarrow\operatorname{Ext}^{i+1}(N,A)
 \qquad(i\ge1).
\end{aligned}
\]
估计 \(B\) 仍只用 \(A,C\) 的同次数项。估计 \(A\) 却要令 \(i>\operatorname{id}B\)、\(i-1>\operatorname{id}C\)，所以其上界是 \(\max\{\operatorname{id}B,\operatorname{id}C+1\}\)；这时左侧的 \(\operatorname{Ext}^{i-1}(N,C)\) 与右侧的 \(\operatorname{Ext}^i(N,B)\) 都为零，正合性使 \(\operatorname{Ext}^i(N,A)=0\)。估计 \(C\) 则要令 \(i>\operatorname{id}B\)、\(i+1>\operatorname{id}A\)，上界便是 \(\max\{\operatorname{id}B,\operatorname{id}A-1\}\)。因此两处偏移与投射情形不同；对所有 \(N\) 使用内射消失判据即得结论。
\end{proof}

这些不等式也给出精确的一步下降。在 \(0\to K\to P\to M\to0\) 中，若 \(P\) 投射且 \(\operatorname{pd}M=d\ge1\) 有限，则 \(P\ne0\)，故中间项的维数固定为 \(\operatorname{pd}P=0\)。两个不同的不等式于是能分别给出 \(K\) 的上界与下界。第二个投射维数不等式给出
\[
 \operatorname{pd}K\le\max\{\operatorname{pd}P,\operatorname{pd}M-1\}
 =\max\{0,d-1\}=d-1.
\]
第三个投射维数不等式则给出
\[
 d=\operatorname{pd}M\le\max\{\operatorname{pd}P,\operatorname{pd}K+1\}
 =\max\{0,\operatorname{pd}K+1\}.
\]
因 \(d\ge1\)，后一式迫使 \(\operatorname{pd}K\ge d-1\)，从而 \(\operatorname{pd}K=d-1\)。特别地，\(d=1\) 时两式分别给出 \(\operatorname{pd}K\le0\) 与 \(\operatorname{pd}K\ge0\)，所以 \(K\) 是非零投射模。若 \(d=0\)，短列分裂，\(K\) 投射或为零；不能强行套用数值 \(d-1=-1\)。

\subsection{交换代数应用：局部投射维数}
\label{b4:subsec:1002:03}

一般环上的消失判据需要对所有测试模作检验。能否只选一个剩余域作测试，取决于消解在这个域上能保留什么信息：若微分仍非零，单个同调群为零只表示该处的核等于像，并不能说明对应自由项为零。

\ProseParagraphBreak
% 跨卷引用目标：b3:prop:minimal-resolution-existence（17.1.4）。
% 跨卷引用目标：b3:prop:minimal-resolution-betti（17.1.5）。
设 \((R,\mathfrak m,k)\) 为交换 Noether 局部环，\(M\ne0\) 为有限生成的 \(R\) 模。取各项均为有限秩的极小自由消解 \(F_\bullet\to M\)，即每个正次微分满足 \(\mathrm{d}_i(F_i)\subseteq\mathfrak mF_{i-1}\)。因为 \(k=R/\mathfrak m\)，微分矩阵与 \(k\) 张量后全部变为零；向 \(k\) 作 Hom 时，\(\mathfrak m\) 的作用也为零。这样的消解由第三卷命题17.1.4构造，命题17.1.5据此给出
\[
 \operatorname{Tor}_i^R(k,M)\simeq k\otimes_R F_i,\qquad
 \operatorname{Ext}_R^i(M,k)\simeq\operatorname{Hom}_R(F_i,k).
\]
这里 \(R\) 交换，\(k\) 同时视为右 \(R\) 模，Tor 的平衡计算使我们可以用第二变量的自由消解 \(F_\bullet\)。两个复形已无非零微分，故每个整数 \(i\ge0\) 的同调就是该项本身，模 \(M\) 的第 \(i\) 个 Betti 数\footnote{贝蒂（Enrico Betti，1823-1892），意大利数学家，研究代数与拓扑。}为
\[
 \beta_i^R(M)=\operatorname{rank}_R F_i
 =\dim_k\operatorname{Tor}_i^R(k,M).
\]
% 跨卷引用目标：b3:thm:nakayama（3.2.3）。
每个 \(F_i\) 有限生成，第三卷定理3.2.3的 Nakayama 引理使 \(k\otimes_RF_i=F_i/\mathfrak mF_i=0\) 等价于 \(F_i=0\)。还须检查零项之后会不会重新出现非零项。若 \(F_i=0\)，则 \(\mathrm{d}_{i+1}:F_{i+1}\to F_i\) 的目标为零，因而整个 \(F_{i+1}\) 都是其核；消解在 \(F_{i+1}\) 处正合，把这个核识别为 \(\mathrm{d}_{i+2}\) 的像；极小性又使该像包含在 \(\mathfrak mF_{i+1}\) 中。因此
\[
F_{i+1}=\ker\mathrm{d}_{i+1}=\operatorname{im}\mathrm{d}_{i+2}
\subseteq\mathfrak mF_{i+1}.
\]
由 \(F_{i+1}\) 有限生成，再用 Nakayama 引理得到 \(F_{i+1}=0\)。递推后所有更高项都为零；一般正合复形在一个零项后仍可接上非零的两项恒等复形，极小性排除了这种情形。

\ProseParagraphBreak
% 跨卷引用目标：b3:thm:local-pd-vanishing（17.2.3）。
对整数 \(d\ge0\)，若 \(\operatorname{Tor}_{d+1}^R(k,M)=0\)，便有 \(F_{d+1}=0\)，后续各项也为零，因而这条极小消解的长度不超过 \(d\)。反过来，长度不超过 \(d\) 的投射消解也是平坦消解，所以 \(\operatorname{pd}_RM\le d\) 使所有高于 \(d\) 次的 Tor 消失，尤其使这个剩余域测试消失。Ext 的计算同样检测 \(F_i\) 是否为零。因此
\[
\begin{aligned}
\operatorname{pd}_RM
&=\sup\{\,i\ge0:\operatorname{Tor}_i^R(k,M)\ne0\,\}\\
&=\sup\{\,i\ge0:\operatorname{Ext}_R^i(M,k)\ne0\,\},
\end{aligned}
\]
这也解释了第三卷定理17.2.3中只检验剩余域的判据。这里的局部性、极小性及各项有限生成性共同保证检测有效。

\ProseParagraphBreak
例如，设 \(k\) 为域。在 \(R=k[x,y]_{(x,y)}\) 上，剩余域的 Koszul 消解为
\[
 0\longrightarrow R\xrightarrow{\binom{-y}{x}}R^2
 \xrightarrow{(x\ \ y)}R\longrightarrow k\longrightarrow0.
\]
它极小，故 \(\operatorname{Tor}_2^R(k,k)=k\ne0\)，同时长度二给出高次消失，从而 \(\operatorname{pd}_Rk=2\)。相反，\cref{b4:exm:infinite-dimension-dual-numbers}的每一步都保留一个自由项，有限生成并不意味着有限投射维数。在有限投射维数情形，第三卷第17.3节的 Auslander--Buchsbaum 公式\footnote{奥斯兰德（Maurice Auslander，1926-1994），美国数学家，研究同调代数与表示论；布赫斯鲍姆（David Alvin Buchsbaum，1929-2021），美国数学家，研究同调代数及交换代数中的自由消解。}进一步将投射维数与模和底环的深度联系起来。

% !TeX root = ../../Book4.tex
% 纲要标题：### 11.3 零调对象
% 纲要标签：b4:sec:1003
% 纲要位置：计划纲要.md，第 3478 行。

\section{零调对象}
\label{b4:sec:1003}

计算一个指定函子时，消解项不必全部内射，但必须对该函子没有正次导出贡献。仅有原列正合并不足够，逐项的零调条件还需证明。

% 纲要内容：
%
% * 相对于函子的零调性
% * 使用零调消解计算导出函子
% * 不能随意用任意正合复形替代消解
%

\subsection{相对于函子的零调性}
\label{b4:subsec:1003:01}

内射消解对每个左正合函子都可用，但有时它的各项过大，实际计算并不方便。若只计算一个指定函子 \(F\)，可以使用对这个函子没有高次导出贡献的对象。

\begin{definition}\label{b4:def:functor-acyclic}
设 \(\mathcal A,\mathcal B\) 为交换范畴，\(\mathcal A\) 有足够多内射对象，\(F:\mathcal A\to\mathcal B\) 为左正合加性函子。若对象 \(Q\in\mathcal A\) 满足 \(R^iF(Q)=0\) 对所有整数 \(i>0\) 成立，则称 \(Q\) 为 \(F\)-\term[lingtiao-duixiang]{零调对象}[acyclic object]，也称 \(F\)-无上同调对象\termidx[wushangtongdiaoduixiang]{无上同调对象}。给定 \(\mathcal A\) 中的对象 \(M\) 及非负增广上链复形
\[
 0\longrightarrow M\longrightarrow A^0\longrightarrow A^1\longrightarrow\cdots
\]
若此序列正合，且每个 \(A^j\)（\(j\ge0\)）都为 \(F\)-零调对象，则称为 \(M\) 的 \(F\)-\term[lingtiao-xiaojie]{零调消解}[acyclic resolution]。
\end{definition}

这里的零调性描述一个对象相对于指定函子 \(F\) 的高次导出行为，只要求正次数的导出函子消失，不要求 \(F(Q)=0\)。\cref{b4:def:acyclic-complex}中的零调复形则描述整条复形内部的同调，要求所有次数的同调对象为零。零调消解要求各项对 \(F\) 零调；去掉增广项后的复形有 \(H^0(A^\bullet)\cong M\)，而施加 \(F\) 后的正次数上同调也可能非零。

\ProseParagraphBreak
内射对象对每个左正合加性函子都是零调对象；对于固定的函子，零调对象却未必内射。例如取 \(F=\operatorname{Hom}_{\mathbb Z}(\mathbb Z/2\mathbb Z,-)\)。对任意交换群 \(Q\)，长度一自由消解 \(0\to\mathbb Z\xrightarrow{\times2}\mathbb Z\to\mathbb Z/2\mathbb Z\to0\) 给出
\[
 R^1F(Q)\cong Q/2Q,\qquad R^iF(Q)=0\quad(i>1).
\]
因此 \(\mathbb Z/3\mathbb Z\) 对 \(F\) 零调，却因不可除而不是内射交换群；不过此时 \(F(\mathbb Z/3\mathbb Z)=0\)。若改取
\[
 Q=(\mathbb Q/\mathbb Z)\oplus\mathbb Z/3\mathbb Z,
\]
则 \(2Q=Q\)，但 \(3Q\ne Q\)，所以它仍对 \(F\) 零调而非内射。同时 \(F(Q)\cong Q[2]\cong\mathbb Z/2\mathbb Z\ne0\)，明确说明原函子的取值不必消失。

\begin{proposition}\label{b4:prop:acyclic-closure}
设 \(\mathcal A,\mathcal B\) 为交换范畴，\(\mathcal A\) 有足够内射对象，\(F:\mathcal A\to\mathcal B\) 为左正合加性函子，且 \(0\to A\to B\to C\to0\) 为 \(\mathcal A\) 中的短正合列。若 \(A,C\) 都 \(F\)-零调，则 \(B\) 也如此；若 \(A,B\) 都 \(F\)-零调，则 \(C\) 也如此。若 \(B,C\) 都 \(F\)-零调，则 \(A\) 零调当且仅当 \(F(B)\to F(C)\) 为满态射。
\end{proposition}
\begin{proof}
前两项分别由长正合列中的
\(R^iF(A)\to R^iF(B)\to R^iF(C)\) 及
\(R^iF(B)\to R^iF(C)\to R^{i+1}F(A)\) 得到，其中 \(i\ge1\)。当 \(B,C\) 都零调时，片段 \(R^{i-1}F(C)\to R^iF(A)\to R^iF(B)\) 已使 \(R^iF(A)=0\) 对 \(i\ge2\) 成立，但低次部分仍为
\[
 0\longrightarrow F(A)\longrightarrow F(B)\longrightarrow F(C)
 \longrightarrow R^1F(A)\longrightarrow0.
\]
所以唯一可能剩下的正次导出是 \(R^1F(A)\cong\operatorname{coker}(F(B)\to F(C))\)。它消失恰好要求 \(F(B)\to F(C)\) 满，这就是第三种情形多出的条件。
\end{proof}

\subsection{零调消解与导出函子的计算}
\label{b4:subsec:1003:02}

\begin{theorem}\label{b4:thm:acyclic-resolution}
设 \(\mathcal A,\mathcal B\) 为交换范畴，\(\mathcal A\) 有足够内射对象，\(F:\mathcal A\to\mathcal B\) 为左正合加性函子。若 \(0\to M\to A^0\to A^1\to\cdots\) 是 \(F\)-零调消解，则对每个 \(n\ge0\)，有典范同构
\[
 R^nF(M)\simeq H^n(F(A^\bullet)).
\]
此同构对增广消解之间的态射自然。
\end{theorem}
\begin{proof}
令 \(K^0=M\)，从增广嵌入 \(\iota^0:M\hookrightarrow A^0\) 起，逐次取 \(K^{j+1}=\operatorname{coker}\iota^j\)。消解正合性把这个商对象识别为 \(\operatorname{im}(\mathrm{d}_A^j)\)，并给出下一项中的嵌入 \(\iota^{j+1}:K^{j+1}\hookrightarrow A^{j+1}\)。这样前两步是
\[
\begin{gathered}
 0\longrightarrow M\longrightarrow A^0\longrightarrow K^1\longrightarrow0,\\
 0\longrightarrow K^1\longrightarrow A^1\longrightarrow K^2\longrightarrow0,
\end{gathered}
\]
一般的一步写成
\[
 0\longrightarrow K^j\xrightarrow{\iota^j} A^j\xrightarrow{p^j} K^{j+1}\longrightarrow0
 \qquad(j\ge0).
\]
先看一次上同调。微分 \(\mathrm{d}_A^1=\iota^2p^1\)，左正合性使 \(F(\iota^2)\) 单，并把 \(\ker F(p^1)\) 识别为 \(F(K^1)\)。因而施加 \(F\) 后的一次余圈对象是 \(F(K^1)\)，一次余边来自 \(F(A^0)\to F(K^1)\)。第一条短正合列的长正合列又因 \(A^0\) 零调而给出 \(F(A^0)\to F(K^1)\to R^1F(M)\to0\)，所以
\[
 H^1(F(A^\bullet))
 \cong\operatorname{coker}\bigl(F(A^0)\to F(K^1)\bigr)
 \cong R^1F(M).
\]
一般次数的计算把同一机制用于后面的短正合列。长正合列和 \(A^j\) 的零调性给出
\[
\begin{gathered}
 0\to F(K^j)\to F(A^j)\to F(K^{j+1})\to R^1F(K^j)\to0,\\
 R^iF(K^{j+1})\simeq R^{i+1}F(K^j)\qquad(i\ge1).
\end{gathered}
\]
消解的微分分解为
\[
 \mathrm{d}_A^n:\quad A^n\xrightarrow{p^n}K^{n+1}
 \xrightarrow{\iota^{n+1}}A^{n+1}.
\]
对第 \(n+1\) 步的短正合列用左正合性，得到 \(F(\iota^{n+1})\) 仍为单态射，所以复合它不改变 \(F(p^n)\) 的核。再对第 \(n\) 步用左正合性，\(F(\iota^n)\) 把 \(F(K^n)\) 识别为这个核。因此
\[
\begin{aligned}
 \ker\bigl(F(\mathrm{d}_A^n)\bigr)
 &=\ker\bigl(F(p^n)\bigr)\\
 &=F(K^n),
\end{aligned}
\]
即 \(F(A^\bullet)\) 在次数 \(n\) 的余圈对象为 \(F(K^n)\)。当 \(n\ge1\) 时，同调就是
\[
 \operatorname{coker}\bigl(F(A^{n-1})\to F(K^n)\bigr)
 \simeq R^1F(K^{n-1})\simeq R^nF(M).
\]
次数零则由左正合性直接等于 \(F(M)\)。这些同构只使用给定的正合列和连接同态，未选分裂，因而对消解态射自然。

为说明这里的“典范”也与最初采用内射消解的定义一致，取 \(M\to I^\bullet\)。下面只使用目标 \(I^j\) 的内射性，源 \(A^j\) 不必内射：先把 \(M\to I^0\) 沿 \(M\hookrightarrow A^0\) 延拓为 \(f^0:A^0\to I^0\)。若已构造到 \(f^j\)，则 \(\mathrm{d}_I^jf^j\) 在 \(K^j\) 上为零；在零次这是与增广相容，在正次这是已有链映射等式及 \(\mathrm{d}_I^j\mathrm{d}_I^{j-1}=0\) 的结果。于是该复合经余核 \(p^j\) 分解为 \(K^{j+1}\to I^{j+1}\)，再由 \(I^{j+1}\) 内射沿 \(\iota^{j+1}\) 延拓为 \(f^{j+1}\)。逐次得到增广上链映射 \(f^\bullet:A^\bullet\to I^\bullet\)：
\[
\begin{tikzcd}[column sep=2.2em,row sep=large]
0\arrow[r]&M\arrow[r]\arrow[d,equal]&A^0\arrow[r]\arrow[d,"f^0"]&A^1\arrow[r]\arrow[d,"f^1"]&\cdots\\
0\arrow[r]&M\arrow[r]&I^0\arrow[r]&I^1\arrow[r]&\cdots
\end{tikzcd}
\]
这里的比较映射通过延拓选取，其同伦唯一性由\cref{b4:prop:comparison-homotopy-unique}的内射版本保证。具体地，设 \(f^\bullet,g^\bullet:A^\bullet\to I^\bullet\) 都延拓 \(1_M\)，约定 \(I^{-1}=0\)、\(h^0=0\)。由于 \(f^0-g^0\) 在 \(M\) 上为零，它经 \(p^0:A^0\twoheadrightarrow K^1\) 分解为 \(K^1\to I^0\)。利用 \(I^0\) 的内射性，把后者沿 \(\iota^1:K^1\hookrightarrow A^1\) 延拓为 \(h^1:A^1\to I^0\)，于是 \(f^0-g^0=h^1\mathrm{d}_A^0\)。

假定已构造到 \(h^j\)（\(j\ge1\)），且低于 \(j\) 次的同伦等式成立。已选的 \(h^j\) 解释了由前一次数传来的差别，在次数 \(j\) 尚未消去的误差是
\[
 e^j=f^j-g^j-\mathrm{d}_I^{j-1}h^j.
\]
链映射等式及前一次的同伦等式给出
\[
\begin{aligned}
 e^j\mathrm{d}_A^{j-1}
 &=\mathrm{d}_I^{j-1}
   \bigl(f^{j-1}-g^{j-1}-h^j\mathrm{d}_A^{j-1}\bigr)\\
 &=\mathrm{d}_I^{j-1}\mathrm{d}_I^{j-2}h^{j-1}=0.
\end{aligned}
\]
这个等式说明误差在前一微分的像上为零。由于 \(\mathrm{d}_A^{j-1}=\iota^jp^{j-1}\) 且 \(p^{j-1}\) 为满态射，\(e^j\iota^j=0\)，故 \(e^j\) 经余核 \(p^j\) 下降为 \(\bar e^j:K^{j+1}\to I^j\)。再由 \(I^j\) 内射，把 \(\bar e^j\) 沿 \(\iota^{j+1}\) 延拓为 \(h^{j+1}:A^{j+1}\to I^j\)。这样逐次得到
\[
 f^j-g^j=\mathrm{d}_I^{j-1}h^j+h^{j+1}\mathrm{d}_A^j
 \qquad(j\ge0),
\]
即 \(f-g=\mathrm{d}_Ih+h\mathrm{d}_A\)。施加加性函子后，这个等式仍成立，所以两种比较映射诱导同一个同调映射。上面的短正合列在比较映射下构成交换图，连接同态自然性说明其诱导映射正是已经构造的同构。因此不同零调消解都通过给定的 \(R^nF(M)\) 典范比较，而不要求它们之间任意方向都有链映射。
\end{proof}

逐项零调并不使 \(F(A^\bullet)\) 成为零调复形。它只使各 \(A^j\) 的正次导出消失，却未要求 \(F(A^{j-1})\to F(K^j)\) 满。上面计算的余核正是这份满射性失败；沿短正合列平移后，它记录的是 \(R^jF(M)\)。

\ProseParagraphBreak
对有足够投射对象的交换范畴及右正合加性函子，对定义域和值域同时取相反范畴，就得到左导出版本：若 \(Q_\bullet\to M\) 正合，且 \(L_iF(Q_j)=0\) 对所有 \(i>0,j\ge0\) 成立，则 \(L_nF(M)\simeq H_n(F(Q_\bullet))\)。在相反范畴中，核与余核互换，内射延拓成为投射提升；再把上链次数改记为链次数，便得到这些正合列与同构，而不需额外假设。具体地，设 \(R\) 为含幺环，\(N\) 为幺性左 \(R\) 模。在幺性右 \(R\) 模范畴中取 \(F=-\otimes_RN\)，平坦右模 \(Q\) 满足 \(L_iF(Q)=\operatorname{Tor}_i^R(Q,N)=0\)（\(i>0\)）。因此任一平坦消解 \(\cdots\to Q_1\to Q_0\to M\to0\) 都给出
\[
 \operatorname{Tor}_n^R(M,N)\cong H_n(Q_\bullet\otimes_RN)
 \qquad(n\ge0),
\]
而不必要求各 \(Q_j\) 投射。

\subsection{正合复形作为消解的条件}
\label{b4:subsec:1003:03}

应用定理时，要检查增广列正合、各项对指定函子零调，并且消解从次数零沿一个方向展开。前两项给出连接同态的平移，最后一项保证每个固定的 \(n\ge2\) 经过有限次就回到 \(M\)：
\[
 R^1F(K^{n-1})\cong R^2F(K^{n-2})\cong\cdots\cong R^nF(K^0)
 =R^nF(M).
\]
一条向两端无限延伸的正合复形若没有这样的增广起点，即使局部短正合列允许平移，也不能据此完成与指定对象导出的比较。

\begin{example}\label{b4:exm:arbitrary-resolution-fails}
取 \(F=\operatorname{Hom}_{\mathbb Z}(\mathbb Z/2\mathbb Z,-)\) 和 \(M=\mathbb Z/2\mathbb Z\)。正合列
\[
 0\longrightarrow M\longrightarrow M\oplus\mathbb Z
 \longrightarrow\mathbb Z\longrightarrow0
\]
用自然包含和投影，确实是一个长度一的右消解。但施加 \(F\) 后，非增广上链复形为 \(\mathbb Z/2\mathbb Z\to0\)，一次上同调为零，而 \(R^1F(M)=\operatorname{Ext}_{\mathbb Z}^1(\mathbb Z/2\mathbb Z,\mathbb Z/2\mathbb Z)=\mathbb Z/2\mathbb Z\)。第零项的零调条件已经失败，因为
\[
 R^1F(M\oplus\mathbb Z)
 \cong (M\oplus\mathbb Z)/2(M\oplus\mathbb Z)
 \cong (\mathbb Z/2\mathbb Z)^{\oplus2}\ne0.
\]
末项 \(\mathbb Z\) 也有 \(R^1F(\mathbb Z)\cong\mathbb Z/2\mathbb Z\ne0\)。原列虽正合，却不是 \(F\)-零调消解，不能用它计算导出函子。
\end{example}

对同一个 \(M=\mathbb Z/2\mathbb Z\)，取
\[
 0\longrightarrow\mathbb Z/2\mathbb Z\xrightarrow{\eta}\mathbb Q/\mathbb Z
 \xrightarrow{\times2}\mathbb Q/\mathbb Z\longrightarrow0,
 \qquad \eta(\bar1)=\dfrac12+\mathbb Z.
\]
乘二映射为满态射，其核恰为 \(\eta(\mathbb Z/2\mathbb Z)\)，而 \(\mathbb Q/\mathbb Z\) 是可除群，因而此列为内射消解。\(F(\mathbb Q/\mathbb Z)\) 由 \(\mathbb Q/\mathbb Z\) 中被 \(2\) 零化的元素确定，同构于 \(\mathbb Z/2\mathbb Z\)；乘二在这些元素上为零，所以施加同一 \(F\) 后，非增广上链复形为
\[
 \mathbb Z/2\mathbb Z\xrightarrow{0}\mathbb Z/2\mathbb Z
 \qquad\text{（次数 }0,1\text{）}.
\]
于是
\[
 R^0F(\mathbb Z/2\mathbb Z)\simeq\mathbb Z/2\mathbb Z,
 \qquad R^1F(\mathbb Z/2\mathbb Z)\simeq\mathbb Z/2\mathbb Z.
\]
这两个正合消解针对同一个对象；前一个含有非零调项，后一个逐项内射，因此后者能够计算该对象的导出函子。

\ProseParagraphBreak
零调对象及消解计算可参阅 Weibel\cite[F-Acyclic Objects 2.4.3, Exercise 2.4.3, Section 2.5]{Weibel1994HomologicalAlgebra}。该书在左导出部分提出此类对象，并以维数平移给出练习；本节把当前所需的右导出版本作为正文定理证明。

% !TeX root = ../../Book4.tex
% 纲要标题：### 11.4 复合函子的导出
% 纲要标签：b4:sec:1004
% 纲要位置：计划纲要.md，第 3484 行。

\section{复合函子的导出}
\label{b4:sec:1004}

把内射消解先送入一个函子以后，其像未必适合计算另一个函子。检查内射像的零调性，才能分清直接复合导出与需要谱序列的情形。

% 纲要内容：
%
% * 函子作用于内射对象后的条件
% * 复合导出的障碍
% * Grothendieck 谱序列的假设预告
%
% ---
%

\subsection{内射对象在函子下的像}
\label{b4:subsec:1004:01}

取一个内射消解后，先作用 \(F\) 再作用 \(G\)，看似能够计算复合函子。但 \(F(I)\) 不一定内射，也不一定对 \(G\) 零调。这一步是复合导出问题中必须单独检查的对象条件。

\begin{proposition}\label{b4:prop:right-adjoint-injectives}
设 \(\mathcal A,\mathcal B\) 为交换范畴，\(U:\mathcal A\to\mathcal B\) 是加性函子，具有正合左伴随 \(V:\mathcal B\to\mathcal A\)。则 \(U\) 把内射对象送到内射对象。
\end{proposition}
\begin{proof}
设 \(I\) 为 \(\mathcal A\) 的内射对象。给定 \(\mathcal B\) 中单态射 \(j:X\hookrightarrow Y\) 及 \(f:X\to U(I)\)，先用伴随双射
\[
 \alpha_X:\operatorname{Hom}_{\mathcal B}(X,U(I))
 \xrightarrow{\sim}\operatorname{Hom}_{\mathcal A}(V(X),I)
\]
把 \(f\) 转置为 \(\widetilde f:V(X)\to I\)。因为 \(V\) 正合，\(V(j):V(X)\hookrightarrow V(Y)\) 仍单；由 \(I\) 内射，可延拓为 \(\widetilde h:V(Y)\to I\)，使 \(\widetilde hV(j)=\widetilde f\)。再令 \(h=\alpha_Y^{-1}(\widetilde h):Y\to U(I)\)。伴随双射对源对象的自然性给出
\[
 \alpha_X(hj)=\alpha_Y(h)V(j)
 =\widetilde hV(j)=\widetilde f=\alpha_X(f).
\]
因 \(\alpha_X\) 为双射，\(hj=f\)，所以所求延拓确实存在，\(U(I)\) 内射。
\end{proof}

例如，设 \(R,S\) 为含幺环，\(\varphi:R\to S\) 为保单位元的环同态，以下模均为幺性左模。限制标量 \(\operatorname{Res}_{\varphi}:S\text{-}\mathbf{Mod}\to R\text{-}\mathbf{Mod}\) 同时出现在两种伴随中：
\[
 S\otimes_R-\;\dashv\;\operatorname{Res}_{\varphi}
 \;\dashv\;\operatorname{Hom}_R({}_RS,-).
\]
张量中的 \(S\) 以 \(t\cdot r=t\varphi(r)\) 接受右 \(R\) 作用；余诱导中的 \({}_RS\) 则以 \(r\cdot t=\varphi(r)t\) 接受左 \(R\) 作用，并在 \(\operatorname{Hom}_R({}_RS,J)\) 上规定左 \(S\) 作用 \((s\cdot h)(t)=h(ts)\)。这些伴随双射已在第二卷命题6.3.4中证明；当前两种保持内射性的判断使用不同的左伴随。

\ProseParagraphBreak
若 \(S\) 作为右 \(R\) 模平坦，则 \(S\otimes_R-\) 正合，所以作为它的右伴随，限制标量保持内射对象。没有平坦性时这可能失败：\(\mathbb Z/2\mathbb Z\) 作为 \(\mathbb F_2\) 模内射，作为整数模却不内射。另一方面，限制标量本身始终正合，因为它保留底交换群和映射；因此它的右伴随，即余诱导 \(\operatorname{Hom}_R({}_RS,-)\)，总把内射左 \(R\) 模送到内射左 \(S\) 模，无需平坦性。

\begin{theorem}\label{b4:thm:derived-composite-exact-first}
设 \(\mathcal A,\mathcal B,\mathcal C\) 为交换范畴，\(\mathcal A,\mathcal B\) 有足够多内射对象。设 \(F:\mathcal A\to\mathcal B\)、\(G:\mathcal B\to\mathcal C\) 为加性函子，\(F\) 正合，\(G\) 左正合，且每个内射对象 \(I\in\mathcal A\) 的像 \(F(I)\) 都对 \(G\) 零调。则对每个对象 \(M\in\mathcal A\) 及整数 \(n\ge0\)，有自然同构
\[
 R^n(GF)(M)\simeq R^nG(F(M)).
\]
\end{theorem}
\begin{proof}
取 \(M\to I^\bullet\) 的内射消解。同一个复形 \(GF(I^\bullet)\) 将用于两种计算。因为 \(I^\bullet\) 是 \(M\) 的内射消解，定义直接给出
\[
 R^n(GF)(M)=H^n(GF(I^\bullet)).
\]
另一方面，\(F\) 正合保证 \(F(M)\to F(I^\bullet)\) 仍为正合消解；内射像对 \(G\) 零调则保证这条消解能够计算 \(G\) 的导出。由\cref{b4:thm:acyclic-resolution}，
\[
 H^n(GF(I^\bullet))\cong R^nG(F(M)).
\]
两条识别复合就是所需同构。前一识别来自导出函子定义，后一识别来自零调消解定理；两者都与消解比较映射相容，所以对 \(M\) 自然。
\end{proof}

这里的内射像条件不能由 \(F\) 的正合性替代。取素数 \(p\)，令 \(F:\mathbb F_p\text{-}\mathbf{Mod}\to\mathbf{Ab}\) 为忘却函子，\(G=\operatorname{Hom}_{\mathbb Z}(\mathbb Z/p\mathbb Z,-)\)。\(F\) 正合，且 \(\mathbb F_p\) 在源向量空间范畴中内射，但
\[
 R^1G(F(\mathbb F_p))
 =\operatorname{Ext}_{\mathbb Z}^1(\mathbb Z/p\mathbb Z,\mathbb Z/p\mathbb Z)
 \cong\mathbb Z/p\mathbb Z\ne0,
\]
所以这个内射对象的像不是 \(G\)-零调对象。同时，在 \(\bar1\) 处取值把 \(GF(V)\) 自然识别为向量空间 \(V\) 的底交换群，故 \(GF\) 也正合，从而
\[
 R^1(GF)(\mathbb F_p)=0
 \quad\text{而}\quad R^1G(F(\mathbb F_p))\cong\mathbb Z/p\mathbb Z.
\]
复合公式在这里确实失败，不能只由前函子正合来保证。

\subsection{复合导出的障碍}
\label{b4:subsec:1004:02}

若 \(F\) 只有左正合性，\(F(I^\bullet)\) 的正次上同调正是 \(R^qF(M)\)，一般不是 \(F(M)\) 的消解。即使其各项都对 \(G\) 零调，也不能在上一个证明中删去 \(F\) 正合这一条件。

\begin{example}\label{b4:exm:composite-first-not-exact}
取素数 \(p\)，令
\[
 F:\mathbf{Ab}\longrightarrow\mathbb F_p\text{-}\mathbf{Mod},\quad F(A)=A[p],
 \qquad G:\mathbb F_p\text{-}\mathbf{Mod}\longrightarrow\mathbf{Ab}
\]
分别为 \(p\)-挠子群函子及忘却函子。\(F\) 左正合，\(G\) 正合，所以每个 \(F(I)\) 都对 \(G\) 零调。可是取 \(M=\mathbb Z\) 的内射消解 \(0\to\mathbb Z\to\mathbb Q\to\mathbb Q/\mathbb Z\to0\)，施加 \(F\) 后，非增广复形为
\[
 F(\mathbb Q)\longrightarrow F(\mathbb Q/\mathbb Z):
 \qquad 0\longrightarrow\mathbb F_p
 \qquad\text{（次数 }0,1\text{）}.
\]
它在一次有非零上同调，因此不是 \(F(\mathbb Z)=0\) 的正合消解。再作用 \(G\) 不会消去这份上同调，故
\[
 R^1(GF)(\mathbb Z)=\operatorname{Ext}_{\mathbb Z}^1(\mathbb Z/p\mathbb Z,\mathbb Z)
 =\mathbb Z/p\mathbb Z,\qquad R^1G(F(\mathbb Z))=0.
\]
非零项来自先作用 \(F\) 时产生的一次导出，而不是来自 \(G\) 的导出。
\end{example}

\begin{proposition}\label{b4:prop:derived-composite-exact-second}
设 \(\mathcal A,\mathcal B,\mathcal C\) 为交换范畴，\(\mathcal A\) 有足够多内射对象，\(F:\mathcal A\to\mathcal B\) 为左正合加性函子，\(G:\mathcal B\to\mathcal C\) 为正合加性函子。则对每个对象 \(M\in\mathcal A\)，有
\[
 R^n(GF)(M)\simeq G(R^nF(M))\qquad(n\ge0)
\]
自然成立。
\end{proposition}
\begin{proof}
这里使用的是 \(G\) 与取同调交换，不需要把 \(F(I^\bullet)\) 看成正合消解。对于任意上链复形 \(C^\bullet\)，正合性使
\[
 Z^n(G(C^\bullet))\cong G(Z^n(C^\bullet)),\qquad
 B^n(G(C^\bullet))\cong G(B^n(C^\bullet)),
\]
因为余圈是微分的核，余边是前一微分的像。再由 \(G\) 保持余核及
\(H^n(C^\bullet)=\operatorname{coker}(B^n(C^\bullet)\hookrightarrow Z^n(C^\bullet))\)，得到
\[
 H^n(G(C^\bullet))\cong G(H^n(C^\bullet)).
\]
取 \(M\to I^\bullet\) 的内射消解，令 \(C^\bullet=F(I^\bullet)\)，便有
\[
 R^n(GF)(M)=H^n(GF(I^\bullet))
 \cong G(H^n(F(I^\bullet)))=G(R^nF(M)).
\]
这些识别来自核、像及余核的自然态射，故与消解比较相容，给出对 \(M\) 的自然性。
\end{proof}

这两个结果处理的是不同的简化情形：前函子正合且把内射对象送到后函子的零调对象时，只剩后函子的导出；后函子正合时，它直接作用于前函子的导出，无需对内射像另加条件。两者都只有左正合性时，两类高次信息必须同时保留。

\subsection{Grothendieck 谱序列的假设}
\label{b4:subsec:1004:03}

设 \(\mathcal A,\mathcal B,\mathcal C\) 为交换范畴，前两者有足够内射对象，\(F:\mathcal A\to\mathcal B\)、\(G:\mathcal B\to\mathcal C\) 均为左正合加性函子。关键附加条件是
\[
 R^pG(F(I))=0\qquad(p>0,\ I\;\text{为}\;\mathcal A\;\text{的内射对象}).
\]
在这些条件下，\secref{b4:sec:1503}中的\cref{b4:thm:grothendieck-spectral-sequence}将用双复形建立 Grothendieck 谱序列，其第二页和目标分别为
\[
 E_2^{p,q}=R^pG(R^qF(M))
 \quad\Longrightarrow\quad R^{p+q}(GF)(M),\qquad p,q\ge0.
\]
指标 \(q\) 记录先作用 \(F\) 时产生的导出次数，\(p\) 记录再对 \(R^qF(M)\) 作用 \(G\) 时的导出次数，最终目标的次数是 \(p+q\)。箭头表示每个目标 \(R^n(GF)(M)\) 带有有限递减滤过
\[
\begin{aligned}
 0&=\mathcal F^{n+1}R^n(GF)(M)\subseteq\mathcal F^nR^n(GF)(M)\subseteq\cdots\\
 &\subseteq\mathcal F^0R^n(GF)(M)=R^n(GF)(M),
\end{aligned}
\]
而极限页描述相邻商：
\[
 E_\infty^{p,n-p}\cong
 \mathcal F^pR^n(GF)(M)/\mathcal F^{p+1}R^n(GF)(M)
 \qquad(0\le p\le n).
\]
因此第二页的各项不能直接读成目标的直和因子；还须经过后续各页，最终得到的也只是滤过的相邻商。谱序列的微分、收敛及滤过将在\chapref{b4:ch:14}建立。上述假设与公式可对照 Weibel\cite[Cohomological Setup 5.8.1, Grothendieck Spectral Sequence Theorem 5.8.3]{Weibel1994HomologicalAlgebra}；该书先作用 \(G\) 再作用 \(F\)，与这里的函子字母次序相反。

\ProseParagraphBreak
在内射像条件成立的前提下，若 \(F\) 正合，则 \(R^qF(M)=0\) 对 \(q>0\) 成立，只剩 \(q=0\) 行；总次数 \(n\) 只可能有 \(E_2^{n,0}=R^nG(F(M))\) 非零。若 \(G\) 正合，内射像条件自动成立，且 \(R^pG(-)=0\) 对 \(p>0\) 成立，只剩 \(p=0\) 列；总次数 \(n\) 只可能有 \(E_2^{0,n}=G(R^nF(M))\) 非零。这两种情形都没有可能的非零后续微分，每个目标滤过也只剩一个相邻商，所以分别恢复前两小节的自然同构。一般情形中，先对 \(F(I^\bullet)\) 作进一步消解，再比较两个取同调次序；内射像的零调条件保证其中一个次序计算 \(R^*(GF)\)。低阶正合列的证明见\cref{b4:prop:spectral-five-term}。

\ProseParagraphBreak
\Cref{b4:tab:derived-composite-conditions}比较这三种情形。表中设 \(\mathcal A,\mathcal B,\mathcal C\) 为交换范畴，\(\mathcal A,\mathcal B\) 有足够多内射对象，\(F:\mathcal A\to\mathcal B\)、\(G:\mathcal B\to\mathcal C\) 为加性左正合函子，\(M\in\mathcal A\)，且 \(n,p,q\ge0\)。

\begin{table}[htbp]
\centering
\caption{复合导出的条件与结论}
\label{b4:tab:derived-composite-conditions}
\begin{tabularx}{\textwidth}{@{}>{\raggedright\arraybackslash}X>{\raggedright\arraybackslash}X@{}}
\toprule
附加条件 & 结论 \\
\midrule
\(F\) 正合，且每个内射对象 \(I\in\mathcal A\) 的像 \(F(I)\) 都对 \(G\) 零调
& \(R^n(GF)(M)\simeq R^nG(F(M))\) \\
\(G\) 正合
& \(R^n(GF)(M)\simeq G(R^nF(M))\) \\
内射对象 \(I\in\mathcal A\) 的像 \(F(I)\) 均对 \(G\) 零调
& \(E_2^{p,q}=R^pG(R^qF(M))\)\newline
\(\Longrightarrow R^{p+q}(GF)(M)\) \\
\bottomrule
\end{tabularx}
\end{table}



\begin{chaptersummary}
有限同调维数可以从消解长度定义，也可以由统一的高次消失刻画。维数平移使一个次数的消失转化为合冲模的次数一判据，因而保证任意消解都能在相应位置截短。局部极小消解另有有限生成和 Nakayama 引理的帮助，所以在那里剩余域足以检测终止性。

\ProseParagraphBreak
零调消解依赖正在计算的函子；任意正合消解并不能替代内射消解。对复合函子 \(GF\)，若 \(F\) 正合且把内射对象送到 \(G\)-零调对象，则只需计算 \(G\) 在 \(F(M)\) 上的导出；若 \(G\) 正合，则可直接将 \(G\) 作用于 \(F\) 的导出。两者都只有左正合性时，在同一内射像零调条件下，Grothendieck 谱序列把前函子的高次导出与后函子的导出共同组织起来。
\end{chaptersummary}

\BookExercises

\begin{exercise}\label{b4:ex:10-dimension-summands}
设 \(R\) 为环，\(A,B\) 为左 \(R\)-模。证明 \(\operatorname{pd}(A\oplus B)=\max\{\operatorname{pd}A,\operatorname{pd}B\}\)，并证明内射维数和平坦维数的同类公式。
\end{exercise}
\begin{solution}
有限直和的消解可逐项直和，给出不超过右侧的上界。反向，\(A,B\) 都是 \(A\oplus B\) 的直和因子；对任意左 \(R\)-模 \(N\)，有
\[
\operatorname{Ext}_R^i(A\oplus B,N)
\simeq\operatorname{Ext}_R^i(A,N)\oplus\operatorname{Ext}_R^i(B,N).
\]
若总模的投射维数至多为 \(d\)，则左端对 \(i>d\) 消失，右侧两项也消失；投射维数判据给出 \(\operatorname{pd}A,\operatorname{pd}B\le d\)。内射维数则要检测第二变量，使用
\[
\operatorname{Ext}_R^i(N,A\oplus B)
\simeq\operatorname{Ext}_R^i(N,A)\oplus\operatorname{Ext}_R^i(N,B).
\]
平坦维数使用任意右 \(R\)-模 \(N\) 及 \(\operatorname{Tor}_i^R(N,A\oplus B)\) 的同类分解。若某个直和项维数无限，总模也不可能具有有限维数；零模约定则与最大值公式相容。
\end{solution}

\begin{exercise}\label{b4:ex:10-flat-not-projective}
固定素数 \(p\)。证明 \(\operatorname{fd}_{\mathbb Z}\mathbb Z[1/p]=0\)、\(\operatorname{pd}_{\mathbb Z}\mathbb Z[1/p]=1\)。
\end{exercise}
\begin{solution}
\(\mathbb Z[1/p]\) 是整数局部化，张量它等于将模中的 \(p\) 变为可逆，正合性可由清除分母直接验证，因此平坦。任意同态 \(f:\mathbb Z[1/p]\to\mathbb Z\) 满足 \(f(1)=p^n\cdot f(1/p^n)\) 对每个 \(n\) 成立，只能有 \(f(1)=0\)，再由整数群无挠得 \(f=0\)。若该模投射，则嵌入一个自由交换群；复合任意坐标投影都为零，嵌入也为零，矛盾。整数整体维数为一，于是其投射维数恰为一。
\end{solution}

\begin{exercise}\label{b4:ex:10-three-base-rings}
设 \(p\) 为素数。分别把 \(\mathbb Z/p\mathbb Z\) 看成 \(\mathbb F_p\)、\(\mathbb Z\)、\(\mathbb Z/p^2\mathbb Z\) 上的模，求其投射维数，并对最后一个底环写出全部正次数 Ext 自群。
\end{exercise}
\begin{solution}
在域上维数为零。整数自由消解 \(0\to\mathbb Z\xrightarrow{p}\mathbb Z\to\mathbb Z/p\mathbb Z\to0\) 给出上界一，\(\operatorname{Ext}_{\mathbb Z}^1(\mathbb Z/p\mathbb Z,\mathbb Z)=\mathbb Z/p\mathbb Z\ne0\) 给出下界一。令 \(S=\mathbb Z/p^2\mathbb Z\)，则
\(\cdots\xrightarrow{p}S\xrightarrow{p}S\to\mathbb Z/p\mathbb Z\to0\)
正合，因为乘 \(p\) 的核与像均为 \(pS\)。对它取 \(\operatorname{Hom}_S(-,\mathbb Z/p\mathbb Z)\) 后全部微分为零，故每个 \(\operatorname{Ext}_S^i(\mathbb Z/p\mathbb Z,\mathbb Z/p\mathbb Z)=\mathbb Z/p\mathbb Z\)，投射维数无限。
\end{solution}

\begin{exercise}\label{b4:ex:10-detect-last-degree}
设 \(R\) 为环，非零左模 \(M\) 的投射维数为有限整数 \(d\)。证明存在左模 \(N\) 使 \(\operatorname{Ext}_R^d(M,N)\ne0\)。内射维数有怎样的对应说法？
\end{exercise}
\begin{solution}
若 \(d=0\)，取 \(N=M\)，恒等映射非零。若 \(d>0\) 且对所有 \(N\) 都有 \(\operatorname{Ext}^d(M,N)=0\)，则消失判据给出 \(\operatorname{pd}M\le d-1\)，与最小性矛盾。内射版本把两变量交换：若 \(\operatorname{id}M=d\)，必有 \(N\) 使 \(\operatorname{Ext}^d(N,M)\ne0\)，证明完全由内射消失判据得到。
\end{solution}

\begin{exercise}\label{b4:ex:10-global-dimension-attained}
设 \(R\) 为非零环，\(J\) 为集合，\((M_j)_{j\in J}\) 为一族左 \(R\)-模。沿用 \(\sup\varnothing=-\infty\) 的约定，证明 \(\operatorname{pd}_R(\bigoplus_jM_j)=\sup_j\operatorname{pd}_RM_j\)。由此证明若左整体维数无限，就存在投射维数无限的单个左模。
\end{exercise}
\begin{solution}
若 \(J\) 为空或所有 \(M_j\) 为零，两边都为 \(-\infty\)。否则，若右侧为有限整数 \(d\ge0\)，为每个 \(M_j\) 选长度至多 \(d\) 的投射消解。逐项直和保持正合，投射模的直和仍投射，故得到长度至多 \(d\) 的消解。反向，每个 \(M_j\) 都是总模的直和因子，前一题的消失判据给出其维数不大于总模维数。因此上确界公式成立，右侧为 \(\infty\) 时也成立。若整体维数无限，或已经有无限维数模，或为每个 \(n\ge0\) 选择 \(\operatorname{pd}M_n>n\)。在后一情形令 \(M=\bigoplus_{n\ge0}M_n\)，则 \(\operatorname{pd}M\ge\operatorname{pd}M_n>n\) 对每个 \(n\) 成立。假如 \(\operatorname{pd}M=d<\infty\)，取 \(n=d\) 就得到 \(d>d\) 的矛盾。因此这些可能各自有限的维数，可以由一个直和合成无限投射维数。
\end{solution}

\begin{exercise}\label{b4:ex:10-syzygy-step}
设 \(R\) 为环，\(0\to K\to P\to M\to0\) 是左 \(R\)-模短正合列，\(P\) 投射。证明 \(\operatorname{pd}M=\infty\) 当且仅当 \(\operatorname{pd}K=\infty\)。当 \(M\) 投射时，\(K\) 是否必须为零？
\end{exercise}
\begin{solution}
如果 \(K\) 有有限长度的投射消解，把它与给定列拼接，就得到 \(M\) 的有限长度投射消解。如果 \(M\) 有有限长度的投射消解，则截短判据或短正合列的不等式使 \(K\) 也有这样的消解；\(M=0\) 时 \(K\cong P\) 投射，结论也直接成立。两者同时不存在有限长度的投射消解，便得到所述无限维数的等价。\(M\) 投射时仅能推出给定列分裂、\(K\) 投射；例如取 \(R=\mathbb Z\)，坐标投影 \(\mathbb Z\oplus\mathbb Z\to\mathbb Z\) 就具有非零核 \(\mathbb Z\)。
\end{solution}

\begin{exercise}\label{b4:ex:10-local-top-tor}
设 \(k\) 为域，\(R=k[x,y]_{(x,y)}\)。利用正文中的剩余域消解求 \(\operatorname{pd}_R(x,y)\)，并计算 \(\operatorname{Tor}_1^R((x,y),k)\)。
\end{exercise}
\begin{solution}
在 \(0\to(x,y)\to R\to k\to0\) 中，\(\operatorname{pd}_Rk=2\)，所以一次合冲模 \((x,y)\) 的投射维数为一。维数平移给出 \(\operatorname{Tor}_1^R((x,y),k)\simeq\operatorname{Tor}_2^R(k,k)=k\)。也可截取消解 \(0\to R\to R^2\to(x,y)\to0\)，张量 \(k\) 后首箭头为零，得到相同结果。
\end{solution}

\begin{exercise}\label{b4:ex:10-acyclic-n-divisible}
设 \(n>1\)，\(F=\operatorname{Hom}_{\mathbb Z}(\mathbb Z/n\mathbb Z,-)\)。证明交换群 \(A\) 对 \(F\) 零调当且仅当 \(nA=A\)。给出一个满足该条件而非内射的群。
\end{exercise}
\begin{solution}
由 \(\mathbb Z/n\mathbb Z\) 的长度一自由消解，\(R^1F(A)=A/nA\)，且 \(R^iF(A)=0\) 对 \(i>1\) 成立，因此零调条件只要求乘这个固定的 \(n\) 满射，即 \(nA=A\)。内射交换群则必须可除，即 \(mA=A\) 对每个非零整数 \(m\) 都成立。取 \(A=\mathbb Z[1/n]\)，它满足前一条件；选择不整除 \(n\) 的素数 \(\ell\)，方程 \(\ell x=1\) 在该群中无解，所以它不满足后一条件，因而不是内射整数模。
\end{solution}

\begin{exercise}\label{b4:ex:10-acyclic-kernel-warning}
对 \(F=\operatorname{Hom}_{\mathbb Z}(\mathbb Z/2\mathbb Z,-)\)，在 \(0\to\mathbb Z\to\mathbb Q\to\mathbb Q/\mathbb Z\to0\) 中找出两个零调项和一个非零调项，并用 \(F\) 后的映射说明区别。
\end{exercise}
\begin{solution}
\(\mathbb Q\)、\(\mathbb Q/\mathbb Z\) 都内射，因而对 \(F\) 零调。\(F(\mathbb Q)=0\)，而 \(F(\mathbb Q/\mathbb Z)\) 由 \(1/2+\mathbb Z\) 生成，是 \(\mathbb Z/2\mathbb Z\)。低次长正合列成为
\[
0\longrightarrow F(\mathbb Q/\mathbb Z)
\xrightarrow{\delta}R^1F(\mathbb Z)\longrightarrow0,
\]
所以连接映射 \(\delta\) 是非零同构 \(\mathbb Z/2\mathbb Z\xrightarrow{\sim}\mathbb Z/2\mathbb Z\)。这正是 \(F(\mathbb Q)\to F(\mathbb Q/\mathbb Z)\) 不满射留下的余核，也是零调对象类不无条件地对满射之核封闭的原因。
\end{solution}

\begin{exercise}\label{b4:ex:10-small-acyclic-resolution}
用 \(0\to\mathbb Z\to\mathbb Z[1/2]\to\mathbb Z[1/2]/\mathbb Z\to0\) 计算 \(F=\operatorname{Hom}_{\mathbb Z}(\mathbb Z/2\mathbb Z,-)\) 在 \(\mathbb Z\) 上的导出函子，并说明这不是内射消解。
\end{exercise}
\begin{solution}
两个非增广项都可被 \(2\) 整除，所以由\cref{b4:ex:10-acyclic-n-divisible}的计算对 \(F\) 零调。\(F(\mathbb Z[1/2])=0\)，而商群中被 \(2\) 零化的元素恰为 \(0\) 和 \(1/2+\mathbb Z\)。因此所得上链复形是 \(0\to\mathbb Z/2\mathbb Z\)，次数为 \(0,1\)，给出 \(R^0F(\mathbb Z)=0\)、\(R^1F(\mathbb Z)=\mathbb Z/2\mathbb Z\)，更高次为零。\(\mathbb Z[1/2]\) 不能被 \(3\) 整除，故不是内射群，却已足以参与计算。所得复形的一次上同调非零：零调消解保证其同调计算 \(R^nF\)，并不保证施加 \(F\) 后仍正合。
\end{solution}

\begin{exercise}\label{b4:ex:10-composite-image-condition}
设 \(p\) 为素数，考虑环同态 \(\mathbb Z\to\mathbb F_p\) 对应的限制标量函子 \(\operatorname{Res}\) 与余诱导函子 \(C=\operatorname{Hom}_{\mathbb Z}(\mathbb F_p,-)\)。取内射交换群 \(I=\mathbb Q/\mathbb Z\)，求 \(C(I)\) 的 \(\mathbb F_p\)-模结构，并计算 \(\operatorname{Ext}_{\mathbb Z}^1(\mathbb Z/p\mathbb Z,\operatorname{Res}C(I))\)。由此判断两个函子各是否保持内射对象，说明这与 \(\operatorname{Res}\) 正合并不矛盾。
\end{exercise}
\begin{solution}
在 \(1\) 处取值把 \(C(I)\) 识别为 \(I[p]\)，后者由 \(1/p+\mathbb Z\) 生成。因此 \(C(I)\simeq\mathbb F_p\) 是一维向量空间，也是内射 \(\mathbb F_p\)-模。限制标量后得到 \(\mathbb Z/p\mathbb Z\)，整数上的长度一自由消解给出
\[
\operatorname{Ext}_{\mathbb Z}^1(\mathbb Z/p\mathbb Z,\operatorname{Res}C(I))
\simeq\operatorname{Ext}_{\mathbb Z}^1(\mathbb Z/p\mathbb Z,\mathbb Z/p\mathbb Z)
\simeq\mathbb Z/p\mathbb Z\ne0.
\]
故 \(\operatorname{Res}\) 不保持内射对象。另一方面，\(\operatorname{Res}\dashv C\)，限制标量始终正合，所以由\cref{b4:prop:right-adjoint-injectives}，余诱导 \(C\) 保持内射对象。\(\operatorname{Res}\) 自身的左伴随却是 \(\mathbb F_p\otimes_{\mathbb Z}-\)，它不正合：张量 \(\mathbb Z\xrightarrow{p}\mathbb Z\) 得到非单射的零映射。这两种伴随对内射性的结论因而不同。
\end{solution}

\begin{exercise}\label{b4:ex:10-coinduction-injective}
设 \(R\to S\) 为环同态，\(I\) 为内射左 \(R\)-模。在 \(\operatorname{Hom}_R(S,I)\) 上规定 \((s\cdot f)(t)=f(ts)\)。证明它是内射左 \(S\)-模，并说明这里不要求 \(S\) 作为 \(R\)-模平坦。
\end{exercise}
\begin{solution}
该公式与左 \(R\)-线性相容，且 \(s_1(s_2f)(t)=f(ts_1s_2)\)，定义左 \(S\)-作用。对左 \(S\)-模 \(M\)，映射在 \(1\) 处取值给出
\(\operatorname{Hom}_S(M,\operatorname{Hom}_R(S,I))\simeq\operatorname{Hom}_R(M,I)\)，逆映射把 \(g\) 送到 \(m\mapsto(s\mapsto g(sm))\)。所以余诱导是限制标量的右伴随，而限制标量始终正合。由\cref{b4:prop:right-adjoint-injectives}，余诱导保持内射；此处左伴随是限制标量，不是张量诱导，因此没有平坦性要求。
\end{solution}

\begin{exercise}\label{b4:ex:10-finite-acyclic-length}
在\cref{b4:thm:acyclic-resolution}的环境中，若 \(M\) 有长度 \(d\) 的 \(F\)-零调消解，证明 \(R^nF(M)=0\) 对 \(n>d\) 成立。为什么这不能推出 \(\operatorname{id}M\le d\)？
\end{exercise}
\begin{solution}
该消解施加 \(F\) 后，在高于 \(d\) 的次数没有项，所以其同调为零，零调消解定理给出所需消失。这只针对一个 \(F\)，而内射维数要求全部 \(\operatorname{Ext}^{d+1}(N,M)\) 消失。例如取 \(F=\operatorname{Id}\)，它正合，故 \(R^iF(M)=0\) 对所有对象 \(M\) 及所有 \(i>0\) 成立。于是任意 \(M\) 都有长度零的零调消解 \(0\to M\xrightarrow{1_M}M\to0\)，对 \(M\) 是否内射没有任何限制；取 \(M=\mathbb Z\) 就不是内射群。因此这种消解的长度只针对所选函子，不能读作对象的内射维数。
\end{solution}

\begin{exercise}\label{b4:ex:10-injective-dimension-integers}
求 \(\mathbb Z\)、\(\mathbb Q\)、\(\mathbb Q/\mathbb Z\) 的整数模内射维数，并说明为什么 \(\mathbb Z\) 的投射维数与内射维数不同。
\end{exercise}
\begin{solution}
\(\mathbb Q\) 与 \(\mathbb Q/\mathbb Z\) 都非零且可除，所以内射维数为零。\(0\to\mathbb Z\to\mathbb Q\to\mathbb Q/\mathbb Z\to0\) 给出 \(\operatorname{id}_{\mathbb Z}\mathbb Z\le1\)，而 \(\operatorname{Ext}_{\mathbb Z}^1(\mathbb Z/2\mathbb Z,\mathbb Z)=\mathbb Z/2\mathbb Z\ne0\) 给出等号。\(\mathbb Z\) 自身自由，所以投射维数为零；两种维数分别限制 Ext 的第一、第二变量，没有一般的相等要求。
\end{solution}

\begin{exercise}\label{b4:ex:10-dimension-bounds-sharpness}
在整数模范畴中，对\cref{b4:prop:dimension-short-exact}的三条投射维数不等式及三条内射维数不等式，分别构造取等和严格不等的短正合列。可以让同一条短正合列同时检验多个不等式。
\end{exercise}
\begin{solution}
固定素数 \(p\)。先取三条短正合列
\[
\begin{aligned}
\xi_1:&\quad0\longrightarrow\mathbb Z\xrightarrow{\,p\,}\mathbb Z
\longrightarrow\mathbb Z/p\mathbb Z\longrightarrow0,\\
\xi_2:&\quad0\longrightarrow\mathbb Z\longrightarrow\mathbb Q
\longrightarrow\mathbb Q/\mathbb Z\longrightarrow0,\\
\xi_3:&\quad0\longrightarrow\mathbb Z/p\mathbb Z\longrightarrow\mathbb Z/p\mathbb Z\oplus\mathbb Z
\longrightarrow\mathbb Z\longrightarrow0.
\end{aligned}
\]
前两列使用通常的乘法、包含和商映射，第三列使用直和的包含与投影。三个非零项的投射维数依次为 \((0,0,1)\)、\((0,1,1)\)、\((1,1,0)\)。这里 \(\mathbb Q/\mathbb Z\) 不是投射整数模：含非零挠元的群不能嵌入自由交换群，而整数整体维数为一给出上界。在 \(\xi_1\) 中，估计 \(B\) 得到 \(0<\max\{0,1\}=1\)，而估计 \(A,C\) 分别得到 \(0=\max\{0,1-1\}\) 和 \(1=\max\{0,0+1\}\)。在 \(\xi_2\) 中，估计 \(A\) 得到 \(0<\max\{1,1-1\}=1\)，而估计 \(B,C\) 的两式都为 \(1=1\)。在 \(\xi_3\) 中，估计 \(C\) 得到 \(0<\max\{1,1+1\}=2\)，而估计 \(A,B\) 的两式都为 \(1=1\)。所以每一列恰有一条所考察的上界严格，另外两式取等。

内射维数方面，\(\xi_2\) 的三项维数为 \((1,0,0)\)。估计 \(B\) 得到 \(0<\max\{1,0\}=1\)，而估计 \(A,C\) 分别得到 \(1=\max\{0,0+1\}\) 和 \(0=\max\{0,1-1\}\)。\(\xi_1\) 的三项维数为 \((1,1,1)\)：\(\mathbb Z/p\mathbb Z\) 的内射维数为一，因为 \(\operatorname{Ext}_{\mathbb Z}^1(\mathbb Z/p\mathbb Z,\mathbb Z/p\mathbb Z)\ne0\)，而整数整体维数为一。此时估计 \(A\) 得到 \(1<\max\{1,1+1\}=2\)，估计 \(B,C\) 则都为 \(1=1\)。最后取分裂列
\[
0\longrightarrow\mathbb Z\longrightarrow\mathbb Z\oplus\mathbb Q
\longrightarrow\mathbb Q\longrightarrow0.
\]
其三项内射维数为 \((1,1,0)\)。估计 \(C\) 得到 \(0<\max\{1,1-1\}=1\)，估计 \(A,B\) 则都为 \(1=1\)。故六条不等式都可以取等，也都可能严格。
\end{solution}
